\section{引言}

安全对齐让指令模型拒绝危险请求，也可能使模型在合法的敏感主题上产生过度拒答。对可离线研究的开放权重模型，研究者常希望把这种行为变化定位到模型内部，而不是通过新一轮监督微调同时改写大量知识与语言能力。Arditi 等人在 13 个、最大 72B 的开源聊天模型上报告了与拒答相关的一维残差方向，并展示了秩一权重编辑的可行性~\cite{arditi2024refusal}。这一发现提供了机械可解释性线索，却没有自动解决模型、层、组件和强度的选择问题。

人工方向消融的主要成本来自超参数耦合。方向可按层独立使用，也可从某个全局层取得；注意力输出和多层感知机输出对行为的影响并不相同；过强编辑虽可能减少模板化拒答，却也可能显著改变正常提示上的分布。一次可用的干预因此需要同时回答两个问题：怎样降低拒答代理，以及怎样限制相对原模型的漂移。把两者压成未经说明的单一分数会隐藏取舍，也使模型导出的选择依据难以审计。

本文分析的自动化系统将方向发现、参数化消融、评分与导出放入同一命令行流程~\cite{automation2026source}。系统先从两组提示计算逐层方向，再让 Optuna 的树结构 Parzen 估计器搜索方向位置和各组件的层权重核；每个试验都在重置的低秩适配器上运行，并同时返回拒答关键词率与首响应词元的 Kullback--Leibler（KL）散度。最终界面展示 Pareto 最优试验，而不是用隐藏权重代替使用者选择。这一工程组织使机制、参数、分值和导出工件能够相互追踪。

本文作出四项具体贡献：

\begin{itemize}
  \item 依据固定提交 \code{bedb94e} 还原系统的真实执行顺序，区分方向发现数据、scorer 评估数据、普通残差均值路径与研究分析路径。
  \item 给出浮点方向插值、组件特定层权重核、投影更新以及预归一化与全量归一化低秩实现的统一数学描述，并标出随机低秩近似的精度边界。
  \item 结合 Qwen3.8-27B 官方模型卡，解释 Gated DeltaNet 与 Gated Attention 输出投影的统一映射，以及默认思考模板对残差和散度测量切点的影响。
  \item 建立公开案例证据卡，把第三方报告的数值、来源 revision、未报告字段和复现限制并列呈现，避免将模型卡结果改写为本文实验。
\end{itemize}

本文不是新的 27B 模型复现实验，也不把拒答降低解释为模型质量提高。IEEE 建议方法应提供足以支持复现的细节，并要求结果讨论承认限制、避免夸大~\cite{ieeestructure}；因此，后文对每项主张分别标记源码事实、官方声明、同行评审证据或第三方报告。通用 \code{IEEEtran} 版式来自 IEEE 提供的会议模板路径~\cite{ieeeauthoring}，但具体会议的语言、页限、匿名和伦理规则仍需在投稿前按当年征稿说明核验。
